\section{Introduction}
Movie services offer more films than an individual can browse. Ratings provide a sparse signal of preference because most viewers rate only a small fraction of a catalog. The resulting task is to estimate ratings for unseen films from observed user--item interactions. The evaluation therefore focuses on rating prediction error; lower error alone does not establish better ranking quality or user satisfaction.

Three approaches are evaluated: a training-set mean, a user--movie bias model, and alternating least squares (ALS) with latent user and movie vectors. The factor model captures interaction patterns shared across the dataset and is widely used in recommender systems \citep{koren2009}. Genre information is evaluated as a prior, with particular attention to items with limited rating history. This design distinguishes standard prediction from cold-start settings for users or movies with sparse observations.

The experiments address four questions: the predictive value of personalization over mean and bias baselines; the accuracy--cost trade-off associated with model capacity; the effect of genre information under natural and controlled cold-start definitions; and the feasibility of serving local recommendations to a new user from a saved model. Root mean squared error (RMSE) measures rating prediction, while recommendation examples illustrate model use without constituting a ranking evaluation.

The saved model is also integrated into a local application that accepts ratings, excludes previously rated films, and presents recommendations. The application is a systems demonstration; clicks, ranked-list quality, and user satisfaction were not measured. MovieLens enables comparison on a large public dataset, but the findings apply to this dataset and split and should not be generalized to all viewers without further evaluation.

The study contributes a comparison of these methods using all 32 million MovieLens ratings and a fixed train/test split; analyses of model capacity, user activity, and natural versus controlled cold start; and a local application that fits a new-user profile without retraining movie parameters. Ten-factor ALS attains a test RMSE of 0.78, compared with 1.06 for the training-mean baseline and 0.86 for the bias model. Genre information modestly improves the controlled cold-start result but not the naturally unrated-item result. Because the analysis uses one random split, these estimates are comparative evidence rather than confirmation of gains on future data.
